% ============================================================
% Cheat sheet, argomento 8: il diagramma dei casi d'uso (dispensa, capitolo 6)
% ============================================================
\argomento{Il diagramma dei casi d'uso}{Dispensa: capitolo 6. Chi usa il sistema e per fare che cosa}[Casi d'uso]

\voce{attore}{Attore}{un ruolo, non una persona}
  {\begin{tikzpicture}
     \ucactor{a1}{(0,0)}{Cliente}
     \ucactor{a2}{(1.9,0)}{Servizio pagamenti}
     \ucactor{a3}{(3.9,0)}{Sensore}
   \end{tikzpicture}}
  {Un'entità \textbf{esterna} che interagisce con il sistema: una classe di utenti con lo stesso ruolo, un altro
   sistema, l'ambiente fisico. Il nome è un \textbf{sostantivo}. La stessa persona può avere più ruoli (più attori).}

\voce{casouso}{Caso d'uso}{ellisse con un verbo}
  {\begin{tikzpicture}
     \node[uc, minimum width=3.2cm] at (0,0) {Prelevare contanti};
   \end{tikzpicture}}
  {Una funzionalità che dà un \textbf{beneficio} a un attore. Il nome contiene un \textbf{verbo}. Un caso d'uso è la
   \textbf{classe} di molti scenari con lo stesso obiettivo (come una classe e i suoi oggetti). Premere un pulsante o
   vedere un errore \textbf{non} sono casi d'uso: sono passi.}

\voce{associazioneuc}{Associazione}{l'attore esegue il caso d'uso}
  {\begin{tikzpicture}
     \ucactor{cu}{(0,0.5)}{Cliente}
     \ucactor{ca}{(0,-0.9)}{Cassiere}
     \node[uc, minimum width=2.0cm] (pc) at (2.7,-0.1) {Pagare};
     \ucactor{bp}{(5.0,-0.1)}{Banca}
     \draw[assoc] ([xshift=9pt]cu) -- (pc);
     \draw[assoc] ([xshift=9pt]ca) -- (pc);
     \draw[assoc] (pc) -- node[pos=0.75, above, font=\tiny]{0..1} ([xshift=-9pt]bp);
   \end{tikzpicture}}
  {Una linea fra attore e caso d'uso. Più attori sulla \textbf{stessa} ellisse \textbf{collaborano} in ogni
   esecuzione (cliente e cassiere insieme), non «uno o l'altro». La cardinalità \code{0..1} dice che la banca
   partecipa solo a volte (pagamento con carta).}

\voce{confine}{Confine del sistema}{rettangolo con il nome del sistema \oltre}
  {\begin{tikzpicture}
     \draw[line width=0.7pt, draw=black!70] (1.1,-1.0) rectangle (4.6,1.0);
     \node[font=\tiny\bfseries, anchor=north west] at (1.15,0.98) {Bancomat};
     \node[uc, minimum width=2.6cm] (p) at (2.85,0.3) {Prelevare};
     \node[uc, minimum width=2.6cm] (s) at (2.85,-0.5) {Vedere il saldo};
     \ucactor{c}{(0,-0.1)}{Cliente}
     \draw[assoc] ([xshift=9pt]c) -- (p);
     \draw[assoc] ([xshift=9pt]c) -- (s);
   \end{tikzpicture}}
  {Il rettangolo racchiude i casi d'uso del sistema: gli attori stanno \textbf{fuori}. Le slide non lo disegnano, ma è
   la notazione standard e aiuta a separare ciò che il sistema fa da chi lo usa.}

\voce{generalizzazioneattori}{Generalizzazione tra attori}{triangolo vuoto verso l'attore generale}
  {\begin{tikzpicture}
     \ucactor{cu}{(0,1.0)}{Cliente}
     \ucactor{rc}{(0,-1.0)}{Cliente abituale}
     \node[uc, minimum width=2.4cm] (bi) at (3.2,1.0) {Comprare};
     \node[uc, minimum width=2.4cm] (fp) at (3.2,-0.9) {Scegliere il premio};
     \draw[assoc] ([xshift=9pt]cu) -- (bi);
     \draw[assoc] ([xshift=9pt]rc) -- (fp);
     \draw[gen] (0,-0.35) -- (0,0.25);
   \end{tikzpicture}}
  {Il cliente abituale \textbf{è un} cliente: può fare tutto ciò che fa il cliente (comprare) e in più ha i suoi casi
   d'uso. È il modo giusto per dire «lo possono fare entrambi». Un attore specializzato senza casi d'uso propri è
   inutile.}

\voce{include}{«include»}{dal caso base a quello incluso: sempre}
  {\begin{tikzpicture}
     \node[uc, minimum width=2.0cm] (wm) at (0,0.55) {Prelevare};
     \node[uc, minimum width=2.0cm] (dc) at (0,-0.55) {Versare};
     \node[uc, minimum width=2.4cm] (ci) at (3.4,0) {Identificare\\la carta};
     \draw[dep] (wm) -- node[stereo, above, sloped]{«include»} (ci);
     \draw[dep] (dc) -- node[stereo, below, sloped]{«include»} (ci);
   \end{tikzpicture}}
  {La parte \textbf{comune} a più casi d'uso diventa un caso a sé, che i casi base eseguono \textbf{ogni volta}. Freccia
   tratteggiata \textbf{dal base all'incluso}: il base dipende dall'incluso (prelevare senza identificare la carta non ha
   senso). Nel codice: una chiamata che c'è sempre.}

\voce{extend}{«extend»}{dall'estensione al caso base: a volte}
  {\begin{tikzpicture}
     \node[uc, minimum width=2.6cm, minimum height=1.05cm] (wm) at (0,0) {\textbf{Prelevare}\\[-2pt]
        \rule{1.8cm}{0.4pt}\\[-2pt]{\tiny \textit{extension points:}}\\[-3pt]{\tiny ricevuta}};
     \node[uc, minimum width=2.0cm] (pr) at (3.6,0) {Stampare\\ricevuta};
     \draw[dep] (pr) -- node[stereo, above]{«extend»} node[stereo, below]{(ricevuta)} (wm);
   \end{tikzpicture}}
  {Un comportamento \textbf{facoltativo} o sotto \textbf{condizione}, aggiunto al caso base in un punto preciso
   (\textit{extension point}, scritto nell'ellisse). Freccia \textbf{dall'estensione al base}: il base ha senso anche
   da solo, l'estensione no. Nel codice: un'aggiunta che il base non conosce.}

\voce{generalizzazioneuc}{Generalizzazione tra casi d'uso}{il caso specializzato sostituisce il generale}
  {\begin{tikzpicture}
     \node[uc, minimum width=2.8cm] (ci) at (1.6,0.6) {\textit{Identificare}};
     \node[uc, minimum width=2.0cm] (g1) at (0,-0.7) {con il PIN};
     \node[uc, minimum width=2.0cm] (g2) at (3.2,-0.7) {con l'impronta};
     \draw[gen] (g1) -- (ci);
     \draw[gen] (g2) -- (ci);
   \end{tikzpicture}}
  {Un caso d'uso che è un \textbf{modo diverso} di fare quello generale, anche molto diverso. Differenza con
   «extend»: l'estensione \textbf{aggiunge} qualcosa al base in un punto preciso, la specializzazione lo
   \textbf{sostituisce}.}

\begin{guida}{Dal testo dell'esercizio al diagramma dei casi d'uso}
\centering
\renewcommand{\arraystretch}{1.3}
\begin{tabular}{@{}l l@{}}
  \textbf{Nel testo} & \textbf{Nel diagramma} \\ \hline
  un sostantivo che indica chi usa il sistema (o un sistema esterno) & attore \\
  un verbo che indica un obiettivo dell'utente (\textit{prenotare}, \textit{pagare}) & caso d'uso \\
  «X e Y insieme», «con l'aiuto di», «tramite il servizio di pagamento» & più attori sulla stessa ellisse \\
  «X oppure Y», «anche gli amministratori possono» & generalizzazione tra attori \\
  «ogni volta», «in entrambi i casi bisogna», passi ripetuti in più casi & «include» \\
  «se vuole», «opzionalmente», «può anche», «in caso di» & «extend» \\
  «in due modi», «con il PIN o con l'impronta» & generalizzazione tra casi d'uso \\
  «dopo aver fatto il login», «se è autenticato» & \textbf{precondizione} nel testo, non «include» \\
\end{tabular}
\end{guida}

\begin{guida}{La specifica di un caso d'uso: il template fully-dressed di Cockburn}
Il diagramma non basta: per ogni caso d'uso si scrive una descrizione testuale. Campi del template (in inglese):
\textbf{Use Case \#} e nome, \textbf{Goal in Context}, \textbf{Preconditions}, \textbf{Success End Condition},
\textbf{Failed End Condition}, \textbf{Primary Actor}, \textbf{Trigger}; poi il \textbf{Main Scenario} (passi numerati,
una colonna per attore e una per il sistema), le \textbf{Extension} (ognuna parte da un passo, con una condizione, e
termina o torna al main scenario) e gli \textbf{Open Issues} (vuoti alla consegna). Formati più leggeri: \textbf{brief}
(un paragrafo) e \textbf{casual} (più paragrafi), per le prime fasi.
\end{guida}

\begin{attenzione}{Errori tipici}
\begin{itemize}[nosep, leftmargin=*]
  \item Casi d'uso che non sono obiettivi dell'utente: «premere OK», «vedere l'errore del PIN» sono \textbf{passi}.
  \item Due attori sulla stessa ellisse per dire «lo può fare l'uno o l'altro»: serve la \textbf{generalizzazione}.
  \item «include» per dire «prima»: il login è una \textbf{precondizione}, non un caso incluso.
  \item Frecce al contrario: «include» va dal base all'incluso, «extend» dall'estensione al base.
  \item Diagrammi pieni di frecce incrociate: i dettagli vanno nella descrizione testuale.
\end{itemize}
\end{attenzione}
